\chapter{Indian Feminism and Environmentalism}

\section{Indian Feminism: The Four Phases}

\begin{KeyConcept}
\begin{itemize}
\item The Indian women's movement is divided into phases: \textbf{social reform during the freedom movement} (1850--1950s), \textbf{1947--1975}, \textbf{1976--1985}, and the continuing phase.
\item \textbf{First phase} --- initiated by men and the British: the \textbf{Abolition of Sati (William Bentinck, 1829)}, the \textbf{Widow Remarriage Act (I. C. Vidyasagar, 1856)}, the \textbf{Female Infanticide Prevention Act}, women's education (Vidyasagar, Raja Ram Mohan Roy, Dayanand Saraswati), and the \textbf{Age of Consent Act (1891, Lord Lansdowne} --- raising the age from 10 to 12). Women-led: \textbf{Savitribai Phule} (first girls' school, 1848), \textbf{Tarabai Shinde} (Stri Purush Tulna), \textbf{Pandita Ramabai}.
\item \textbf{Second phase (1915--1960)} --- the nationalist movement overtook women's concerns; \textbf{Gandhi} initiated women into the movement using feminine methods (sacrifice, fasting, tolerance). Women's organisations: the \textbf{All India Women's Conference} (an NGO founded by \textbf{Margaret Cousins}, affiliated with the Indian National Congress), the \textbf{National Federation of Indian Women} (the women's wing of the CPI), the \textbf{Women's India Association (1917)} and the \textbf{National Council of Women in India (1925)}.
\item \textbf{Post-1947} --- nation-building took precedence; the Constitution gave franchise and civic rights (\textbf{Article 14}, DPSPs --- equal pay for equal work, maternal/child care). The \textbf{state adopted a patronizing role} (women as a 'weaker section'), so \textbf{Indian women did not have to struggle for basic rights as Western women did} (who fought against a patriarchal state).
\end{itemize}
\end{KeyConcept}

\begin{KeyConcept}
\begin{itemize}
\item From the \textbf{1970s}: \textbf{welfare politics $\rightarrow$ protest politics}. Feminists challenged \textbf{unequal wages, relegation to unskilled work, and women's use as a 'reserve army of labour'} --- the aim was to \textbf{abolish the free service of women} (women as cheap capital). Feminist class-consciousness spread across tribal, Dalit, Muslim and lower-class women. \textbf{1975} was the International Decade for Women.
\item \textbf{Left-party-based women's organisations} (SEWA in Gujarat, the Working Women's Forum in Tamil Nadu, the Shramik Mahila Sangathan in Maharashtra) worked for the \textbf{economic upheaval} of women --- fighting \textbf{public patriarchy}. \textbf{Autonomous women's organisations} (no party/trade-union affiliation) worked for \textbf{political and social upheaval} --- politicising violence against women, and fighting \textbf{private patriarchy} (dowry, domestic violence). They \textbf{shunned the welfare approach and adopted protest politics}.
\item Results: the \textbf{anti-dowry movement} (Dowry Prohibition Amendment Act, 1984 --- only limiting, not banning, dowry), the \textbf{anti-sati movement} (after the \textbf{Roop Kanwar} case --- the Commission of Sati Prevention Bill), and the \textbf{Chipko movement}.
\end{itemize}
\end{KeyConcept}

\section{The Two Goals: Corporate and Liberal Feminism}

\begin{KeyConcept}
\begin{itemize}
\item The Indian women's movement worked for \textbf{two goals} (\textbf{Jana Everett's classification}):
\item \textbf{1. Liberation/upliftment of women} --- reforming social practices so women can play a constructive role --- called \textbf{corporate feminism}.
\item \textbf{2. Equal rights for men and women} --- extending to women the civil rights men enjoy (political, economic, familial) --- called \textbf{liberal feminism} (which works through law, asking the state for equal representation --- not overthrowing patriarchy).
\item Other classifications: \textbf{Kalpana Shah} (moderate/women's-rights, radical and socialist feminism); \textbf{Sangari \& Vaid} (modernising patriarchal modes of regulating women vs democratising gender relations).
\end{itemize}
\end{KeyConcept}

\section{The Issues: Brahmanism and the Dalit-Woman Question}

\begin{CriticalPoint}
\begin{itemize}
\item \textbf{The big issue of the women's movement}: the \textbf{concerns of low-caste women are completely absent}. Both the left-party and autonomous organisations (inspired by Western feminism) \textbf{did not address Brahmanism and Brahmanical patriarchy} --- the specifically Indian issues of controlling female sexuality.
\item \textbf{All women do not suffer the same kind or level of discrimination}: upper-caste women face dowry deaths and family control over mobility/sexuality; \textbf{Dalit women face the collective, public threat of rape, sexual assault and physical violence}. Most rape in India is against low-caste women, and the \textbf{conviction rate for rape of Dalit women is under 2\%} (against \textasciitilde{}25\% overall) --- because of the caste nexus between the police, the perpetrator and the political class (highlighted by Kancha Ilaiah). (Nirbhaya and the Hyderabad victim were upper-caste; the Hathras Dalit victim's body was burnt by the administration.)
\item The problem: \textbf{'all Dalits are assumed to be males, and all women are assumed to be a uniform category'} --- leading to the \textbf{universalization of the upper-caste woman's experience (and the Dalit male's experience)}.
\item \textbf{Christine Delphy}: the home too is a \textbf{mode of production} --- the upper-caste housewife employs a maid (a Dalit woman) as her 'clerk', and exercises her power over her.
\end{itemize}
\end{CriticalPoint}

\section{Dalit Feminism and the Autonomous Dalit Women's Organisations}

\begin{KeyConcept}
\begin{itemize}
\item In the \textbf{1990s}, \textbf{Autonomous Dalit Women's Organisations} emerged (the \textbf{National Federation of Dalit Women, the All India Dalit Women's Forum, the Dalit Mahila Samiti}) --- underlining the \textbf{Brahmanism/casteism of the mainstream feminist movement and the patriarchy of Dalit politics}. The split in Indian feminism thus came in the form of \textbf{upper-caste vs Dalit women's organisations} (like the West's liberal/radical split).
\item \textbf{Dalit women suffer triple oppression}: (1) oppressed by the upper castes (casteism), (2) by the Hindu patriarchal system (control of sexuality), and (3) by \textbf{Dalit men} (patriarchy within). They are pushed into bonded labour, the devadasi system, sex trafficking, witch-hunting and prostitution.
\item \textbf{Sharmila Rege}: Dalit feminism challenges the \textbf{'masculinization of Dalithood'} (the Dalit movement as a Dalit-men's movement) and the \textbf{'savarnization of womanhood'} (the women's movement as an upper-caste women's movement). The \textbf{Dalit feminist standpoint is emancipatory} --- it interrogates the Brahminical middle-class biases of earlier feminists.
\item \textbf{SRED} (Society for Rural Education and Development, a Tamil Nadu NGO, founded 1979): recognising that \textbf{land ownership} is central to fighting caste-based discrimination (Bina Agarwal), it formed women's collectives that identified and occupied wastelands --- giving Dalit women land ownership, dignity and status.
\end{itemize}
\end{KeyConcept}

\begin{ExampleBox}
Contemporary issues: \textbf{son meta-preference} (noted by the Economic Survey --- more among the dependent upper-caste woman), \textbf{sex-selective abortion} (sex-determination tests are legally banned but continue), the \textbf{declining sex ratio}, and \textbf{sex trafficking} (largely a lower-caste issue --- see the International Justice Mission).
\end{ExampleBox}

\section{Gail Omvedt: What Counts as a Women's Movement}

\begin{KeyConcept}
\begin{itemize}
\item \textbf{Gail Omvedt} (an American sociologist who settled in India in 1983 and took Indian citizenship) distinguished four kinds of movement --- only the last is a true women's movement:
\item 1. \textbf{Movements where women participate} --- general movements (anti-price-rise, nationalist) in which women are present but the family structure is not the focus --- participation, not action.
\item 2. \textbf{Movements of women} --- movements on general issues (slum improvement, price rise) where women are the only participants --- but this \textbf{may confirm the gender division of labour} (men fight for wage rise, women against price rise) and is not a women's movement.
\item 3. \textbf{Women's reform movements} --- education, abolition of sati (pre-independence) --- which did \textbf{not challenge the fundamental structure of oppression} in family and society (so dowry could simply be redefined as 'gift').
\item 4. \textbf{Women's liberation movements} --- the real women's movement, channeled by the ideology of \textbf{fighting the sexual division of labour and patriarchy} (e.g. the demand for \textbf{'wage for housework'} in Tamil Nadu politics).
\item The \textbf{Chipko movement} is only a 'movement of women' (a fight for livelihood), not a women's movement --- the women did not question the sexual division of labour that made them collect forest produce.
\end{itemize}
\end{KeyConcept}

\section{Ecological Feminism: Val Plumwood and Vandana Shiva}

\begin{KeyConcept}
\begin{itemize}
\item \textbf{Val Plumwood} --- the Western rationalist worldview runs on \textbf{dualism} (human/nature, reason/emotion, mind/body, masculinity/femininity): \textbf{masculinity is identified with mind, logic and reason; women with emotion, body and nature}. This dualistic framework \textbf{sanctions the domination of both women and nature}. Hence \textbf{'ecology is feminine and development is patriarchal'} --- the same force exploits women and nature.
\item \textbf{Vandana Shiva} (the Indian eco-feminist): \textbf{'the modern idea of development is based on Western patriarchy'} --- science, rationality and development are a \textbf{Western patriarchal project}, responsible for the \textbf{subjugation of women and the destruction of nature}.
\item In ancient times, \textbf{nature was seen as a living being} and identified with the \textbf{nurturing mother} --- \textbf{prakriti as a feminine principle} (nature provides food and is fertile, just as the mother does). People were \textbf{accountable to nature} (the Ramayana and Mahabharata show nature being revered and permission being sought). Land was \textbf{communal} --- the earth was a goddess; the idea of \textbf{private property over nature is a Western patriarchal idea}.
\item With modern science, \textbf{prakriti is viewed as dead, separate from man, inferior --- and exploited for machines and finished goods}. The human--nature relation shifted from \textbf{symmetrical/horizontal to vertical} (man controls nature). The reaction to this subjugation: the \textbf{Chipko movement, Narmada Bachao, the Silent Valley (Kerala) and Kudankulam (Tamil Nadu) protests}.
\item This is the answer to the PYQ: \textbf{'Western patriarchy which surrenders feminine principles is the new development project in India --- comment.'}
\end{itemize}
\end{KeyConcept}

\section{The Evolution of Western Environmentalism}

\begin{KeyConcept}
\begin{itemize}
\item Environmental consciousness developed \textbf{late in the West} (India had it in its ethics). Its stages:
\item \textbf{1. The conservation era (mid-1800s)} --- inspired by \textbf{Henry David Thoreau}: preserve the environment; environmental laws to regulate industry and protect wilderness; groups like the \textbf{Sierra Club} and the \textbf{Earth Liberation Front}; national parks by legislation (\textbf{Yellowstone, 1872}; the Grand Canyon).
\item \textbf{2. The modern environmental movement (1950 onwards)} --- a response to ecological disasters: \textbf{Rachel Carson's 'Silent Spring'} criticised \textbf{DDT} (the US EPA later banned it) and brought environmental issues to the public, ushering in \textbf{'green democracy'}. (In India, a parallel movement arose against the pesticide \textbf{Endosulfan}.)
\item \textbf{3. Mainstream environmentalism (1970 onwards)} --- the \textbf{Sierra Club and Greenpeace became institutionalised pressure groups}, providing scientific and economic expertise and acquiring land for preservation (forums like Rio+20, the IPCC).
\item \textbf{4. Grassroots environmentalism (1980s)} --- bottom-up: local volunteers and tribal groups, drawing on \textbf{feminism, Marxism (ecological Marxism) and spirituality}, and cutting across ethnic, racial, gender and class lines.
\item America even has a \textbf{Green Party} (the only party whose agenda is environmentalism --- in a two-party system explained by \textbf{Duverger's law}).
\end{itemize}
\end{KeyConcept}

\begin{QuickRevision}
\begin{itemize}
\item Indian feminism: social reform (men/British) $\rightarrow$ nationalist phase (Gandhi's feminine methods; AIWC/NFIW) $\rightarrow$ 1970s protest politics (left-party orgs = public patriarchy/economics; autonomous orgs = private patriarchy/dowry/domestic violence; anti-dowry 1984, anti-sati/Roop Kanwar, Chipko).
\item Two goals: corporate feminism (liberation/upliftment) + liberal feminism (equal rights) --- Jana Everett; also Kalpana Shah; Sangari \& Vaid.
\item Issues: Brahmanism \& Brahmanical patriarchy ignored; Dalit women's concerns absent (rape conviction <2\% vs \textasciitilde{}25\%); 'all Dalits = male, all women = uniform'; Christine Delphy (home as mode of production).
\item 1990s: autonomous Dalit women's organisations; Sharmila Rege ('masculinization of Dalithood, savarnization of womanhood'); triple oppression; SRED (land ownership); Dalit feminist standpoint emancipatory.
\item Gail Omvedt: movements where women participate / of women / women's reform / women's liberation (only the last is a true women's movement; Chipko is 'of women').
\item Eco-feminism: Val Plumwood (dualism --- masculinity=mind/reason, women=nature/emotion; 'ecology is feminine, development is patriarchal'); Vandana Shiva (development = Western patriarchal project; prakriti as nurturing mother $\rightarrow$ dead/exploited; Chipko/Narmada/Silent Valley/Kudankulam).
\item Western environmentalism: conservation (Thoreau, Yellowstone, Sierra Club) $\rightarrow$ modern (Rachel Carson's Silent Spring, DDT) $\rightarrow$ mainstream (Greenpeace institutionalised) $\rightarrow$ grassroots (feminism/Marxism/spirituality); Green Party.
\end{itemize}
\end{QuickRevision}

